\documentclass{article}
\usepackage[UTF8]{ctex}
\usepackage{amsmath,amssymb,amsfonts}
\usepackage{enumitem}
\usepackage{geometry}
\geometry{a4paper, margin=1in}
% === 错题本排版优化（不参与题目标识） ===
\setlength{\parindent}{0pt}
\setlength{\parskip}{0.35em plus 0.10em minus 0.08em}
\setlength{\emergencystretch}{3em}
\setlist{itemsep=0.25em, topsep=0.40em, parsep=0pt}
\allowdisplaybreaks[2]
% === 错题本排版优化结束 ===
\begin{document}
\section*{上海交通大学2026年数学分析真题}
\section*{44．上海交通大学2026年数学分析试卷}

\subsection{题目1}
1．计算下列极限

\[
\lim _{x \rightarrow 0} \frac{1+\cos x-\cos x \cos (2 x)-\sqrt[3]{\cos x}}{\ln (1-x)+\ln (1+x)}
\]

\subsection{题目2}
2．设 \(u(x, y)\) 在 \(\mathbb{R}^{2}\) 上具有二阶连续偏导数，并且满足

\[
\frac{\partial^{2} u}{\partial x^{2}}-\frac{\partial^{2} u}{\partial y^{2}}=0
\]

若 \(u(x, 2 x)=x, u_{x}^{\prime}(x, 2 x)=x^{2}\) ，求 \(u_{x x}^{\prime \prime}(x, 2 x), u_{x y}^{\prime \prime}(x, 2 x), u_{y y}^{\prime \prime}(x, 2 x)\) ．

\subsection{题目3}
3．设 \(f(x, y, z)\) 是 \(\Omega=\left\{(x, y, z) \in \mathbb{R}^{3} \mid x^{2}+y^{2}+z^{2} \leq 1\right\}\) 上的黎曼可积函数，\(p \in \mathbb{R}\) ，若 \(|f(x, y, z)|\)在原点附近有正下界，试求广义积分

\[
I=\iiint_{\Omega} \frac{f(x, y, z) \ln \sqrt{x^{2}+y^{2}+z^{2}}}{\left(x^{2}+y^{2}+z^{2}\right)^{p}} \mathrm{~d} x \mathrm{~d} y \mathrm{~d} z
\]

收敛时 \(p\) 的取值范围．

\subsection{题目4}
4．设 \(f(x)\) 在区间 \([0, a]\) 上有二阶连续导数，\(f^{\prime}(0)=1, f^{\prime \prime}(0) \neq 0\) ，且 \(0<f(x)<x, x \in(0, a)\) ，任取\(x_{1} \in(0, a)\) ，令 \(x_{n+1}=f\left(x_{n}\right), n=1,2, \cdots\) ，证明：数列 \(\left\{x_{n}\right\}\) 和 \(\left\{n x_{n}\right\}\) 都收敛．

\subsection{题目5}
5．证明：若函数 \(f(x)\) 在 \(\mathbb{R}\) 上有定义，并且满足：\(\forall x \in \mathbb{R}, \exists \delta_{x}>0\) ，当 \(x_{1}, x_{2} \in\left(x-\delta_{x}, x+\delta_{x}\right)\) 且\(x_{1}<x<x_{2}\) 时，必有 \(f\left(x_{1}\right)>f(x)>f\left(x_{2}\right)\) ，则 \(f(x)\) 在 \(\mathbb{R}\) 上是严格单调减少函数．

\subsection{题目6}
6．设 \(f(x)\) 是 \((0,+\infty)\) 上的一致连续函数，\(\left\{\lambda_{n}\right\}\) 是正数列，若对任意固定的 \(x \in(0,+\infty)\) ，有

\[
\lim _{n \rightarrow \infty} f\left(x+\lambda_{n}\right)=0
\]

证明：函数列 \(\left\{f\left(x+\lambda_{n}\right)\right\}\) 在闭区间 \([a, b](a>0)\) 上一致收敛于零．

\subsection{题目7}
7．设 \(f(x)\) 是 \(\mathbb{R}\) 上有界的连续函数，若存在正常数 \(a\) ，使得

\[
f(x)+a \int_{x-1}^{x} f(t) \mathrm{d} t
\]

为 \(\mathbb{R}\) 上的常值函数，证明：\(f(x)\) 是 \(\mathbb{R}\) 上的常值函数．
\end{document}
